\section{Discussion}

% Briefly discuss your choices of methodology, (technical) solutions other results.

The choice of a RAG architecture over a fine-tuned model reflects a pragmatic trade-off. RAG allows grounding in verified curriculum content without the cost and data requirements of fine-tuning, and makes the system's knowledge base transparent and editable by educators. However, it limits the depth of pedagogical reasoning the system can perform — the quality of scaffolding depends heavily on what is in the retrieval corpus and how well the MiniGuide prompt constraint shapes the LLM's output.

The decision to use Google Gemini as the LLM introduces a dependency on a third-party, non-EU cloud service. This is acceptable for prototyping but would need to be resolved before any classroom deployment. The planned migration to a local CNN for handwriting recognition and an EU-based text-only LLM addresses this, though it introduces a different trade-off: a smaller, local model may produce lower-quality scaffolding than a frontier multimodal model.

The evaluation methodology is limited in scope. Two sessions — one with teachers, one with students — can validate the interaction design and surface usability issues, but cannot measure learning outcomes, long-term engagement, or the effectiveness of the adaptive difficulty system. These would require longitudinal studies with larger participant groups.